Fig.~\ref{fig:nominal} shows both transitions on the identified plant for every controller and Table~\ref{tab:nominal} the metrics. HOT reaches cruise in \hotFrontT{}~s and the hover-exit airspeed in \hotBackT{}~s with the lift fraction never below \hotFrontLf{} and \hotBackLf{}, the angle of attack below \hotFrontAl{}$^\circ$ and \hotBackAl{}$^\circ$, and \hotFrontEnergy{}~Wh and \hotBackEnergy{}~Wh of energy; the safety module intervened at \hotFrontInterv{} of the \hotFrontUpd{} front updates and \hotBackInterv{} of the \hotBackUpd{} back updates, and the same policy without the module, HOT$-$LSM, flies the same times (\hotnfFrontT{} and \hotnfBackT{}~s) but crosses the pitch box on the way. The threshold rule is fast in the front transition because it cuts the rotors as soon as the airspeed crosses its threshold, at the price of the lowest lift fraction of all controllers, \thrFrontLf{}; the schedule and the gain-scheduled PID keep the lift near unity but need \schFrontT{}~s and \pidFrontT{}~s in the front and 11 to 15~s in the back transition; the nominal-model MPC keeps the lift fraction at unity and takes \mpcFrontT{}~s with its \SI{4}{s} horizon and rate limits. The feed-forward and recurrent policies, run with the same module, take \mlpFrontT{}~s and \lstmFrontT{}~s in the front transition and hold less lift; SAC reaches the target in \sacFrontT{}~s and \sacBackT{}~s, and PPO does not reach cruise within the window on the identified plant, its rotor command oscillating between its bounds (Fig.~\ref{fig:nominal}d).
\begin{figure}[t]\centering\includegraphics[width=\columnwidth]{fig06_nominal.pdf}
\caption{Nominal front (left) and back (right) transitions: airspeed, lift fraction with the \SI{75}{\percent} bound, angle of attack with its bound, rotor thrust, for eight controllers (SAC, close to PPO, is in Table~\ref{tab:nominal}).}\label{fig:nominal}
\end{figure}
\begin{table*}[t]\centering\footnotesize
\caption{Metrics on the identified plant.}\label{tab:nominal}
\begin{tabular}{l|cccc|cccc}\toprule
& \multicolumn{4}{c|}{front transition ($6\to15$ m/s)} & \multicolumn{4}{c}{back transition ($15\to6$ m/s)}\\
controller & time [s] & energy [Wh] & min $\lf$ & max $\alpha$ [deg] & time [s] & energy [Wh] & min $\lf$ & max $\alpha$ [deg]\\\midrule
THR & 4.6 & 2.99 & 0.84 & -0.0 & 12.2 & 3.51 & 0.95 & 3.2\\
SCH & 7.2 & 3.38 & 0.99 & -0.0 & 11.3 & 3.77 & 0.98 & 0.8\\
GS-PID & 5.2 & 3.32 & 0.99 & -0.0 & 14.5 & 3.56 & 0.98 & -0.0\\
MPC & 18.4 & 3.39 & 1.00 & 0.4 & 19.5 & 3.60 & 1.00 & 0.3\\
MLP & 3.8 & 4.02 & 0.92 & -0.1 & 6.9 & 4.51 & 0.92 & 1.9\\
LSTM & 6.6 & 4.16 & 0.97 & -0.0 & 5.2 & 4.87 & 0.94 & -0.8\\
HOT$-$LSM & 4.0 & 4.17 & 0.94 & -0.2 & 4.7 & 4.51 & 0.96 & 0.6\\
PPO & -- & 8.94 & 0.32 & 1.2 & 8.2 & 4.03 & 0.94 & 3.2\\
SAC & 7.3 & 3.01 & 0.91 & -0.3 & 4.9 & 4.30 & 0.94 & 4.7\\
\textbf{HOT} & 4.5 & 4.21 & 0.94 & -0.2 & 4.7 & 4.51 & 0.96 & 0.6\\
\bottomrule\end{tabular}
\end{table*}
